\section{What the measurement cannot see}
\label{sec:symmetry}

Let $\pi$ be the prior density and assume the likelihood depends on $x$ only through $\cA(x)$, as for \cref{M-eq:operator}.
Since every $g\in\cG$ is a permutation, it preserves Lebesgue measure, and \cref{lem:orbit} gives at once:

\begin{proposition}
\label{prop:ratio}
For all $x$ and $g\in\cG$: $\pi(x)\,p(gx\mid y)=\pi(gx)\,p(x\mid y)$. Where $\pi(x)>0$ and $p(x\mid y)>0$, this reads $p(gx\mid y)/p(x\mid y)=\pi(gx)/\pi(x)$.
\end{proposition}

Within an orbit, the posterior therefore ranks images by the prior alone.

\begin{corollary}[measurement-based selection is orientation-blind]
\label{cor:blind}
Let $S(\hat x_1,\dots,\hat x_K;y)$ select one of $K$ runs using only $\cA(\hat x_1),\dots,\cA(\hat x_K)$ and $y$; the residual rule $\Eo$ is an example.
Then $S$ returns the same index when any run is replaced by any member of its orbit.
The same holds for weights or resampling probabilities that depend on a particle only through its likelihood $p(y\mid\hat x)$, such as residual-based pruning~\cite{makris2026tgd} and likelihood-only potentials in sequential Monte Carlo.
Full twisted-SMC weights~\cite{wu2023tds} also contain prior and proposal factors, which need not be invariant.
\end{corollary}

\Cref{cor:blind} needs no assumption on the prior, the noise or the solver\claim{C-T2}.
It explains why the measurement residual can remove runs that do not fit $y$ but can never correct a rotated run, and why particle methods that prune by prior consistency rather than likelihood reduce orientation errors~\cite{kocabas2026pcdaps}.
It concerns the measurement, not the quality of the runs: a solver may produce rotated runs that fit $y$ worse than its upright runs, and for two of our three solvers it does (\cref{M-sec:selection}).

\paragraph{Where the preference for upright images comes from.}
Consider a sampler with Gaussian transitions, $x_{i-1}=m_i(x_i)+s_i\xi_i$ with $\xi_i\sim\mathcal N(0,I)$ and state-independent $s_i>0$, started from a $\cG$-invariant law; $y$ is fixed and dropped from the notation.
Deterministic steps, such as ODE denoising, are absorbed into the next mean map.
DAPS and SITCOM-ODE have this form up to the Langevin iterate they return; FAST-DIPS returns a deterministic denoised estimate of its last state, so the results below cover that state, not its output\claim{C-T3}.
Measurement gradients, projections in measurement space, isotropic noise and pixel-wise clipping commute with $\cG$; the denoiser need not.

\begin{theorem}[equivariant samplers]
\label{thm:equivariant}
If every mean map is $\cG$-equivariant, $m_i(gx)=g\,m_i(x)$, the law of the output $x_0$ is $\cG$-invariant.
If, moreover, the eight members of the orbit of $\xs$ are distinct and $B$ is a neighborhood of $\xs$ whose images $gB$ are pairwise disjoint, a run that lands in $\bigcup_g gB$ lands in each $gB$ with probability $1/8$.
\end{theorem}

The argument is the invariance of marginals under equivariant kernels used in equivariant generative models~\cite{xu2022geodiff,hoogeboom2022edm}, applied to an inverse problem with a symmetric forward operator.
It is the diffusion counterpart of the classical observation that iterative algorithms reconstruct an object or its twin with equal probability~\cite{fienup1986stagnation}, and of the finding that guidance cannot select between modes its gradient does not distinguish~\cite{dandapanthula2026tilting}.
A learned denoiser is not equivariant: an upside-down face is not a typical face.
The next result bounds how much this matters.

\begin{theorem}[orientation preference is bounded by the equivariance defect]
\label{thm:tv}
Let $\Delta^g_i(x)=m_i(x)-g\,m_i(g^{-1}x)$ and let $L$ be the law of $x_0$.
For every $g\in\cG$ and every set $B$,
\begin{align}
\bigl|\Prob(x_0\in B)-\Prob(x_0\in gB)\bigr| &\le \TV(L,g_\#L)\notag\\
&\le\Bigl(\sum_i\frac{\E\|\Delta^g_i(x_i)\|^2}{4s_i^2}\Bigr)^{1/2}.
\label{eq:tv}
\end{align}
\end{theorem}

The proof compares the chain with its conjugate by $g$, whose output law is $g_\#L$, and combines Pinsker's inequality with the chain rule for the KL divergence of Gaussian transitions (supplement)\claim{C-T4}.
It requires neither an exact score nor log-concavity nor uniqueness, but it does require Gaussian transitions with positive, state-independent scales up to the returned state.
Taking $B$ as a neighborhood of $\xs$, the left side of \cref{eq:tv} is the gap between the rates of \good runs and of orbit errors through $g$; the bound turns an observed orientation preference into a lower bound on the sampler's lack of equivariance along its own path.
A preference for upright images is not a defect.
If the posterior law $\Pi$ itself prefers upright images, then $\TV(\Pi,g_\#\Pi)\le2\,\TV(\Pi,L)+\TV(L,g_\#L)$ shows that a sampler with an invariant output law stays at least $\TV(\Pi,g_\#\Pi)/2$ away from it (supplement).

\paragraph{Measurements.}
We measured the defect of the denoisers themselves on noised ground-truth images, $x_t=x+\sigma_t\epsilon$ at nine noise levels ($t=100,200,\dots,900$ of the $1{,}000$-step schedule; two draws; the $180$ FFHQ and $100$ ImageNet images).
Relative to the denoising error, $\tilde\delta(t)=\overline{\|D(R_bx_t)-R_bD(x_t)\|^2}\,/\,\overline{\|D(x_t)-x\|^2}$ is $0.10$--$0.31$ on FFHQ but only $0.015$--$0.052$ on ImageNet for the global rotation $R_{111}$, against $0.26$--$0.61$ and $0.15$--$0.43$ for a single-channel rotation $R_{100}$ (\cref{tab:s2-c2}).
In this diagnostic, the normalized global-flip defect of the ImageNet denoiser is much smaller than that of the FFHQ denoiser.
Single-channel rotations produce larger defects on both datasets.
This is a proxy: \cref{thm:tv} involves the defect of the whole transition along the sampler's own trajectory.
On FFHQ, $59$ of the $1{,}960$ runs of the three solvers are orbit errors ($58$ global, $1$ partial): $7.5\%$ of the failed runs, from $3.4\%$ for FAST-DIPS to $13.0\%$ for DAPS-2K\claim{C-A1}.
On ImageNet the share is far higher: $1{,}105$ of the $6{,}400$ FAST-DIPS runs on ImageNet are orbit errors ($1{,}047$ global, $58$ partial), $22.6\%$ of the failed runs; it is $22.4\%$ on ImageNet$'$ and $31.0\%$ for DAPS-4K ($163$ of its $525$ failed runs).
These observations are qualitatively consistent with denoiser-dependent orientation preferences.
They do not test \cref{thm:tv} directly.
Its bound concerns the sampler's transition defects along its own trajectories, not denoiser defects on noised ground-truth images.

\begin{table}[h]
\centering\small
\setlength{\tabcolsep}{4pt}
\begin{tabular}{@{}lccccccccc@{}}
\toprule
$t$ & 100 & 200 & 300 & 400 & 500 & 600 & 700 & 800 & 900 \\
\midrule
FFHQ, $R_{111}$ & .101 & .150 & .192 & .229 & .257 & .274 & .305 & .274 & .219 \\
FFHQ, $R_{100}$ & .563 & .597 & .611 & .604 & .580 & .555 & .542 & .440 & .264 \\
ImageNet, $R_{111}$ & .026 & .035 & .042 & .050 & .052 & .048 & .029 & .015 & .025 \\
ImageNet, $R_{100}$ & .427 & .399 & .380 & .364 & .336 & .297 & .238 & .189 & .150 \\
\bottomrule
\end{tabular}
\caption{\textbf{Relative equivariance defect $\tilde\delta(t)$ of the denoisers on noised ground-truth images.} Ratio of means over images and two noise draws; the bounds of the image-bootstrap $95\%$ intervals lie within $7\%$ of each value for FFHQ and within $28\%$ for ImageNet.}
\label{tab:s2-c2}
\end{table}

\paragraph{Correcting orientation without the ground truth.}
By \cref{prop:ratio}, the Bayes rule within an orbit maximizes $\pi$, so any test with the same ordering of orbit members is admissible.
On FFHQ, a zero-shot CLIP~\cite{radford2021clip} test with fixed prompts (``a photo of a face.'', ``a photo of a person's face.'', ``a portrait photo of a person.'', and the same three with ``an upside-down'' in place of ``a'') fixes all $58$ global orbit errors among the $1{,}960$ runs and flips none of the $1{,}171$ \good runs, at $8.5$\,ms per candidate.
A score built from the solvers' shared denoiser at intermediate noise levels does the same and, since it compares all eight members, also fixes the one partial orbit error\claim{C-A2}.
A generic real-vs-fake detector~\cite{ojha2023univfd}, used as an orientation score, is the least reliable test on ground-truth images (\cref{tab:s2-gt}).
Orientation correction is therefore generic and cheap on faces.
On ImageNet, where orientation is less canonical, the generic prompts ``a photo.'' and ``an upside-down photo.'' repair $57$ of the $67$ global orbit errors among FAST-DIPS's first four runs and flip $5$ of $100$ \good runs (\cref{M-sec:protocol}).
We also applied every test to the ground-truth images themselves (\cref{tab:s2-gt}).
On FFHQ, all tests but one keep all $180$ images upright.
On ImageNet, the generic prompts keep $92$ of $100$ upright, and the denoiser scores keep $93$ to $95$.
The errors fall mostly on different images.
No ImageNet image is turned by all of the generic CLIP test, the two global-pair denoiser scores and UnivFD, and $20$ are turned by at least one.

\begin{table}[h]
\centering\small
\setlength{\tabcolsep}{4pt}
\begin{tabular}{@{}lcc@{}}
\toprule
Test & FFHQ ($180$) & ImageNet ($100$) \\
\midrule
CLIP, face prompts & 180 & --- \\
CLIP, generic prompts & 180 & 92 \\
CLIP, generic prompt ensemble & 180 & 92 \\
Denoiser score $S_{\mathrm{mid}}$, global pair & 180 & 95 \\
Denoiser score $S_{\mathrm{mid}}$, all eight members & 180 & 95 \\
Denoiser score at one level ($t=300$), global pair & 180 & 93 \\
UnivFD, all eight members & 172 & 89 \\
\bottomrule
\end{tabular}
\caption{\textbf{Orientation tests on the ground-truth images:} how many each keeps upright. Over all eight members, $S_{\mathrm{mid}}$ makes the same choice as over the global pair on every image. One of UnivFD's errors is a channel-wise flip.}
\label{tab:s2-gt}
\end{table}
Restricted to $\hat x$ and $R_{111}\hat x$, the denoiser score matches the generic CLIP test on ImageNet and leaves $5$ points less orientation loss on ImageNet$'$ (\cref{tab:s4-prior}; exploratory).
